\section{Springer actions and the supported Hall Hopf structure}
\label{sec:supported-springer}

A proper flag remembers an ordering of the successive line quotients
even when the middle sheaf does not split. Over distinct support points,
permuting those quotients gives the ordinary permutation action on the
finite flag cover. Smallness extends that action across collisions.
Its trace will identify the Hall multiplication with a shuffle product
on the coefficient objects of Section~\ref{sec:support-coalgebra}.
The same objects already carry the actual direct-sum coproduct.

\subsection{The convolution tensor and full flags}

Retain the rational Betti category, common base $B_E$, and objects
$\mathcal B_n,\mathcal T_n,\mathcal V_n$ of
Section~\ref{sec:support-coalgebra}. For coefficients on $Y_a,Y_b$,
write
\begin{equation}\label{sh:convolution}
 \mathcal C_a*\mathcal D_b
 =Rm_{a,b,*}(\mathcal C_a\boxtimes_{B_E}\mathcal D_b).
\end{equation}
Finite support addition and the external-product comparisons give
the associativity identifications used below. The unit is
$\mathbb Q_{B_E}$ on $Y_0$. In particular, the ordered line
cochains have coherent identifications
\[
 \mathcal V_a*\mathcal V_b\simeq\mathcal V_{a+b},
 \qquad \mathcal V_0=\mathbb Q_{B_E}.
\]
Let
\begin{equation}\label{sh:ordered-inclusion}
 j_n:\mathcal B_n\longrightarrow\mathcal V_n
\end{equation}
be $\theta_n$ from \eqref{sc:theta} followed by the invariant
inclusion. Theorem~\ref{sc:ordered-equivalence} makes $j_n$ a
split monomorphism. At this point its invariant description uses
the ordinary permutation of the ordered supports and line factors.

Let $\overline{\mathcal F}_n$ be the stack of full flags
\[
 0=F_0\subset F_1\subset\cdots\subset F_n,
 \qquad \operatorname{length}(F_i)=i,
\]
on one common elliptic torsor. There are maps
\[
 h_n:\overline{\mathcal F}_n\longrightarrow
                    \overline{\mathcal M}_n,
 \qquad
 g_n:\overline{\mathcal F}_n\longrightarrow\widetilde L_n.
\]
The first forgets the flag; the second records
$(F_1/F_0,\ldots,F_n/F_{n-1})$ in this order. Thus the convention
is subobject first, as in the Hall product. Put
\[
 \mathrm{Spr}_n=Rh_{n,*}\mathbb Q,
 \qquad
 u_n:\mathbb Q\longrightarrow\mathrm{Spr}_n,
 \qquad
 t_n:\mathrm{Spr}_n\longrightarrow\mathbb Q.
\]
Here $u_n$ is the adjunction unit and $t_n$ is the proper
complex-oriented full-flag trace with the normalized coefficient
$+1$ of Proposition~\ref{gh:coefficient-maps}. These are maps on
$\overline{\mathcal M}_n$; no support pushforward has yet been
taken. For $n=0$ both maps are the identity on $B_E$.

\subsection{Smallness and the trace identities}

\begin{lemma}\label{sh:smallness}
The map $h_n$ is proper and small. Both its source and target are
smooth stacks of complex dimension $-1$. Over the distinct-support
open substack it is the cover of ordered flags, with group $S_n$.
\end{lemma}
\begin{proof}
The matrix neighbourhoods of Lemma~\ref{sc:neighbourhoods}
identify a full flag with a complete invariant flag of the matrix.
Before quotienting by a change of frame, the local map is a
restriction of
\[
 \widetilde{\mathfrak g}
 =\{(A,F_\bullet):AF_i\subset F_i\text{ for all }i\}
 \longrightarrow\mathfrak g=\mathfrak{gl}_n(\mathbb C).
\]
The source is a vector bundle over the complete flag variety,
smooth of dimension $n^2$, and the map is proper. A matrix with
distinct eigenvalues has precisely its $n!$ ordered eigenline
flags. The curve version of these matrix charts and the flag
cover is also described in \cite[\S\S3.1--3.3]{FengYunZhang2024}.

We check smallness at the collision strata. Let $\mathfrak b$
be the upper-triangular matrices. For a pair of complete flags of
relative position $w\in S_n$, the common preserving matrices
have dimension $\dim\mathfrak b-\ell(w)$. The orbit of that pair
has dimension $\dim\operatorname{Flag}_n+\ell(w)$. Thus each
relative-position stratum in
$\widetilde{\mathfrak g}\times_{\mathfrak g}
\widetilde{\mathfrak g}$ has dimension $n^2$.
In its common preserving matrix space, matrices with distinct
eigenvalues form a nonempty dense open set: a Cartan diagonal in
a common adapted basis contains such matrices. The complement
of the regular semisimple locus in each of these finitely many
strata therefore has dimension at most $n^2-1$.

For a stratum $Z$ of $\mathfrak g$ outside that locus, let $d_Z$
be the dimension of the flag fibre, constant after refining the
stratification. Its inverse-image fibre product has dimension
$\dim Z+2d_Z$. The preceding bound gives
\[
 \dim Z+2d_Z<n^2,
\]
which is the smallness inequality. It persists under restriction
to spectral coordinate neighbourhoods. Changes of frame and
translation preserve these maps and their fibre dimensions, so
the calculation descends to the curve stacks and their common
torsor quotient. The matrix quotient has dimension zero; the
further quotient by the one-dimensional curve $E$ has dimension
$-1$. Smoothness descends from the same charts.
\end{proof}

The small-map theorem identifies the shifted direct image from a
smooth space with the intermediate extension of its covering local
system \cite[\S4.1.14, Remark~9.29, Proposition~9.44]
{deCataldoMigliorini2007}. Apply it on the smooth atlases with
their dimension shifts and descend. In the present common-torsor
convention the perverse object is $\mathrm{Spr}_n[-1]$.
The constant object $\mathbb Q[-1]$ is the intermediate extension
of the trivial summand. This argument uses complex Betti
intermediate extension; the finite-field coefficients of
\cite{FengYunZhang2024} are not needed for the conclusion.

\begin{proposition}\label{sh:springer-trace}
The deck action extends to a coherent action
$\rho_n:S_n\to\operatorname{Aut}(\mathrm{Spr}_n)$, with
\begin{equation}\label{sh:trace-identities}
 t_nu_n=n!\operatorname{id},\qquad
 u_nt_n=\sum_{w\in S_n}\rho_n(w).
\end{equation}
The action convention is inverse pullback by a deck permutation,
so that $\rho_n$ is a left action.
\end{proposition}
\begin{proof}
An intermediate extension has no subobject or quotient supported
on the boundary. A morphism between two such extensions is
therefore determined by its restriction to the open stratum.
The deck operators extend by functoriality of intermediate
extension, and their group equations continue to hold.
On the cover, $u_n$ is the constant tuple and $t_n$ sums its
$n!$ entries. The coefficient normalization and complex
orientation fix this positive sign. Both equations in
\eqref{sh:trace-identities} hold there, hence on the whole stack.

These equations occur in the perverse heart after shifting by
$[-1]$. Its mapping spaces have no positive homotopy groups:
for $k>0$, the corresponding group is
$\operatorname{Hom}(P,Q[-k])=0$. The action and its equalities
thus give a coherent action in the derived enhancement.
Applying a derived functor retains those diagrams.
\end{proof}

\subsection{Associated graded and refinement of blocks}

The endpoint map $g_n$ is not proper. The coefficient comparison
needed from it nevertheless follows from its actual extension
charts.

\begin{lemma}\label{sh:associated-graded}
For the full or any partial flag, associated-graded pullback
induces an equivalence
$\mathcal C\to Rg_*g^*\mathcal C$ for the coefficient complexes
pulled from its endpoint stack in the bounded-below rational
category used here. These equivalences are natural in the
coefficient and in base change of the endpoint family.
In particular, there is a canonical equivalence
\begin{equation}\label{sh:gamma}
 \gamma_n:\mathcal V_n\xrightarrow{\ \sim\ }
 \mathcal W_n:=R(\overline c_n h_n)_*\mathbb Q,
 \qquad
 \gamma_n j_n=R\overline c_{n,*}u_n.
\end{equation}
\end{lemma}
\begin{proof}
For a partial flag, the matrix chart consists of block
upper-triangular matrices modulo the preserving parabolic $P$.
The endpoint chart consists of the diagonal blocks modulo its
Levi quotient $L$. After choosing an endpoint family, the
off-diagonal entries form an affine vector bundle and the kernel
of $P\to L$ is a unipotent group bundle. Both are contractible
in the analytic topology. More precisely, each level of the
quotient nerve is an affine bundle over the same endpoint base.
On a local product neighbourhood, cochains with a pulled-back
coefficient identify with the cochains of that coefficient on
the base. This identifies the adjunction unit, not just the
dimensions of its cohomology. The contractions are natural under
restriction and change of the endpoint family. Descent and
iteration prove the stated equivalences. The stabilizer groups
have been included in the quotient nerve; they have not been
discarded from the extension stack.

Since $\overline c_nh_n=q_nr_ng_n$, the full-flag equivalence
gives $\gamma_n$. Its inverse on cochains is restriction to the
ordered split-flag section of $g_n$. Composing that section with
$h_n$ is the actual ordered direct sum $\sigma_n$.
Consequently $\gamma_n^{-1}R\overline c_{n,*}u_n$ is exactly
the adjunction map $j_n$ in \eqref{sh:ordered-inclusion}.
This proves both parts of \eqref{sh:gamma} without a general
nonproper base-change assertion.
\end{proof}

Transport the action through this equivalence:
\[
 \mathcal R_n(w)=\gamma_n^{-1}
          (R\overline c_{n,*}\rho_n(w))\gamma_n
       \quad\text{on }\mathcal V_n.
\]
To obtain a crossing for convolution, these actions must agree
with the actions on blocks.

\begin{proposition}\label{sh:block-action}
For every finite list $\alpha=(n_1,\ldots,n_r)$ with
$\sum_i n_i=n$, the identification
$\mathcal V_{n_1}*\cdots*\mathcal V_{n_r}\simeq\mathcal V_n$
intertwines the convolution of the $\mathcal R_{n_i}$ with
$\mathcal R_n|_{\prod_iS_{n_i}}$. These identifications are
coherent for refinements and zero blocks.
\end{proposition}
\begin{proof}
Let $\overline{\mathcal P}_\alpha$ be the partial-flag stack
with successive quotient lengths $n_i$. Forgetting the steps
inside its blocks gives an actual Cartesian square
\[
 \begin{array}{ccc}
 \overline{\mathcal F}_n&\longrightarrow&
             \displaystyle\prod_{i,B_E}\overline{\mathcal F}_{n_i}\\
 \downarrow&&\downarrow\prod_i h_{n_i}\\
 \overline{\mathcal P}_\alpha&\xrightarrow{\ p_\alpha\ }&
             \displaystyle\prod_{i,B_E}\overline{\mathcal M}_{n_i}.
 \end{array}
\]
Indeed, a refinement is a full flag of each actual successive
quotient. The vertical maps are proper. Proper base change
identifies the direct image along the left map with
$p_\alpha^*(\boxtimes_i\mathrm{Spr}_{n_i})$.
It carries the product of the block Springer actions. Pushing
to $\overline{\mathcal M}_n$ gives endomorphisms of
$\mathrm{Spr}_n$ which, on the distinct-support cover, are the
restrictions of the $S_n$ deck operators. Uniqueness of morphisms
between the same intermediate extensions identifies them
everywhere. This is the proper parabolic refinement mechanism
of \cite[proof of Corollary~3.7]{FengYunZhang2024}, applied here
to the Betti complexes just constructed.

Push to support and use Lemma~\ref{sh:associated-graded} for the
partial flag and its pulled-back block coefficients. The result
is the asserted convolution comparison. Its external-product
identification uses the finite maps $q_{n_i}$ and the constant
line cochains: proper base change and the degreewise K\"unneth
equivalence identify their products with the ordered object
$\mathcal V_n$.

Every further refinement has this same Cartesian diagram of
actual quotient flags. Its action comparisons agree on the
open cover and extend uniquely in the perverse heart. The
derived functors preserve those coherent diagrams. A zero block
inserts an identity map. This proves all finite refinements,
not only an equality of binary cohomology operators.
\end{proof}

In particular, put $\mathcal R=\mathcal R_2((12))$.
Its insertions in $\mathcal V_n$ are the adjacent operators
$\mathcal R_n((i,i+1))$. They satisfy
\[
 \mathcal R_i^2=1,\qquad
 \mathcal R_i\mathcal R_{i+1}\mathcal R_i
  =\mathcal R_{i+1}\mathcal R_i\mathcal R_{i+1},\qquad
 \mathcal R_i\mathcal R_j=\mathcal R_j\mathcal R_i
       \quad(|i-j|>1).
\]
Thus the crossing is defined on the supported cochain objects,
before global cohomology or a polynomial presentation.

\subsection{The fixed Hall product and direct-sum coproduct}

Denote the supported Hall product \eqref{sc:hall-product} by
$\mu_{a,b}:\mathcal B_a*\mathcal B_b\to\mathcal B_{a+b}$.
Its full charge-one product is
\[
 P_n=(R\overline c_{n,*}t_n)\gamma_n:
                           \mathcal V_n\longrightarrow\mathcal B_n.
\]
The coefficient flag equations identify this composite with the
actual iterated Hall product, including nonsplit extensions.

\begin{theorem}\label{sh:operation-comparison}
For every $n\geq0$,
\begin{equation}\label{sh:symmetrizer}
 j_nP_n=\sum_{w\in S_n}\mathcal R_n(w),\qquad
 P_nj_n=n!\operatorname{id}_{\mathcal B_n}.
\end{equation}
Consequently $j_n(\mathcal B_n)$ is the image of the Springer
averaging idempotent. For $a+b=n$, let
$\operatorname{Sh}(a,b)$ consist of the permutations whose images
are increasing within each input block. Then
\begin{align}
 j_n\mu_{a,b}
   &=\sum_{w\in\operatorname{Sh}(a,b)}
           \mathcal R_n(w)(j_a*j_b),\label{sh:shuffle}\\
 (j_a*j_b)\Delta^{\mathrm{sup}}_{a,b}
   &=j_n.\label{sh:coproduct}
\end{align}
These equations compare the fixed geometric operations in the
same ordered object $\mathcal V_n$ and commute with all finite
refinements.
\end{theorem}
\begin{proof}
Push \eqref{sh:trace-identities} to support and use
\eqref{sh:gamma}. This gives \eqref{sh:symmetrizer} and shows
that the Springer invariant image is exactly $j_n(\mathcal B_n)$.
This image equals the ordinary invariant image constructed in
Theorem~\ref{sc:ordered-equivalence}. It does not identify the
two group actions or their complementary averaging projectors.

Full-flag refinement gives
$P_n=\mu_{a,b}(P_a*P_b)$ under
$\mathcal V_n=\mathcal V_a*\mathcal V_b$. Restricting to
$j_a*j_b$ and using \eqref{sh:symmetrizer} gives
\[
 a!b!\,j_n\mu_{a,b}
       =\sum_{w\in S_n}\mathcal R_n(w)(j_a*j_b).
\]
Every permutation is uniquely a shuffle followed by an element
of $S_a\times S_b$. Proposition~\ref{sh:block-action} makes
the latter action trivial on $j_a*j_b$. Cancelling $a!b!$
proves \eqref{sh:shuffle} with its stated order and sign.

Equation~\eqref{sh:coproduct} is the actual ordered direct-sum
comparison in the proof of Theorem~\ref{sc:coalgebra}. The
block invariant images for the Springer action are the same
images with the same inclusions $j_a*j_b$. That inclusion is
split monic, so the equation identifies
$\Delta^{\mathrm{sup}}_{a,b}$ uniquely with the inclusion from
full invariants to block invariants. No transported replacement
of the coproduct, factorial factor, or Euler correction is made.
All equations come from the full flag and direct-sum diagrams
already compared, so they preserve their refinement coherences.
\end{proof}

\subsection{A tensor recipient and the global crossing}

We specify a recipient on which the crossing is natural. Let
$\mathcal P$ be the $\mathbb Q$-linear monoidal category with
objects $n\geq0$, endomorphism algebra $\mathbb Q[S_n]$ at $n$,
and zero morphisms between distinct integers. Tensor concatenates
blocks and uses their group-algebra embeddings. The block
permutation is its symmetry. Take finite sums and idempotent
completion, and allow charge-graded families with the finite
convolution sum in each total charge. Denote this category by
$\mathcal P^{\mathrm{gr}}$.

Proposition~\ref{sh:block-action} gives a coherent monoidal
realization $\mathfrak R$ in the charge-graded support convolution
category, with
\[
 \mathfrak R(n)=\mathcal V_n,
 \qquad \mathfrak R(w)=\mathcal R_n(w).
\]
For the idempotent
$e_n=(n!)^{-1}\sum_{w\in S_n}w$, put $\mathsf S_n=(n,e_n)$.
Equation~\eqref{sh:symmetrizer} identifies
$\mathfrak R(\mathsf S_n)=\mathcal B_n$ with inclusion $j_n$.
No faithfulness of this realization is required. In particular,
we have used the universal group algebras, not their image
quotients defining the global category $\mathcal T_R$.
Nor is $\mathfrak R$ asserted to preserve the ambient ordinary
exchange of arbitrary sheaf coefficients.

\begin{theorem}[Supported Hall Hopf comparison]\label{sh:hopf}
The charge-graded object with components $\mathsf S_n$ is a
connected Hopf object in $\mathcal P^{\mathrm{gr}}$ whose
product is the shuffle sum and whose coproduct is the inclusion
of full invariants into block invariants. Its realization gives
the actual maps $\mu_{a,b}$ and
$\Delta^{\mathrm{sup}}_{a,b}$.
The realized symmetry is the block Springer crossing, the unit
is $\mathcal B_0=\mathbb Q_{B_E}$, the counit is projection to
charge zero, and the antipode on $\mathcal B_n$ is
$(-1)^n\operatorname{id}$.
These operations satisfy all finite-charge Hopf and refinement
identities in this recipient.
\end{theorem}
\begin{proof}
In $\mathbb Q[S_n]$, shuffle sum followed by the block
idempotent has image in $e_n$, and
$(e_a\otimes e_b)e_n=e_n$. Thus both operations are morphisms
between the indicated idempotent objects. The interleaving
argument of Theorem~\ref{gh:formal-hopf} is an identity of
permutations before any representation or image quotient is
taken, so it applies in $\mathcal P^{\mathrm{gr}}$ itself.

More explicitly, an $(a,b)$-interleaving with an output cut after
$p$ positions has unique numbers $i,j$ of entries from the two
input lists before that cut, with $i+j=p$. It splits the inputs
into four ordered lists of lengths $i,a-i,j,b-j$. Exchanging
the middle two lists, and interleaving separately on each side
of the cut, recovers precisely the original permutation.
This is the bijection proving the four-block identity
\eqref{gh:four-block}, now with the universal block permutation.
Further cuts give the same ordered partitions, proving all
regrouping identities. Three input lists similarly prove
associativity; invariant inclusions give coassociativity.

On $\mathsf S_n$ the product after a cut at $a$ is
$\binom na\operatorname{id}$. Hence either convolution of
$(-1)^n\operatorname{id}$ with the identity has coefficient
$\sum_{a=0}^n(-1)^a\binom na$, zero for $n>0$ and one at
$n=0$. This proves both antipode equations. Each fixed charge
uses finitely many cuts; no completion or infinite output sum
is involved. Finally Theorem~\ref{sh:operation-comparison}
identifies the realized product and coproduct, while
Proposition~\ref{sh:block-action} identifies their crossings
and refinements.
\end{proof}

For example, the crossing between
$\mathcal B_a*\mathcal B_b$ and
$\mathcal B_b*\mathcal B_a$ is specified by applying the block
permutation to $j_a*j_b$ and then restricting to the image
$j_b*j_a$. Conjugation carries the first block subgroup to the
second, so this restriction is well defined and invertible.
This describes an actual supported map. Naturality is asserted
for the morphisms of the specified recipient, not for every
map in the ambient category of sheaves.

\begin{corollary}\label{sh:global-crossing}
Under common-base pushforward and the simultaneous comparison
of Theorem~\ref{gh:simultaneous}, the supported crossing has
global image
\[
 R=\sigma_E\sigma_z-\delta_{12}\partial_z.
\]
The supported operations recover
Theorem~\ref{gh:global-hopf}, including its sign and unit.
\end{corollary}
\begin{proof}
The global coproduct comparison is
Proposition~\ref{sc:global-comparison}. The proper flag trace
and its endpoint map push forward to the same global Hall
product. In charge two, \eqref{sh:symmetrizer} reads
$j_2P_2=1+\mathcal R$. The global calculation identifies this
same geometric composite with
$1+\sigma_E\sigma_z-\delta_{12}\partial_z$.
Subtracting the identity proves the formula. Block compatibility
then identifies all higher crossings and their operation
diagrams.
\end{proof}

\begin{example}\label{sh:binary-example}
The two flags over distinct support points produce the two
terms in $1+\mathcal R$. Their extension across a collision is
fixed by the proper flag complex, not by an arbitrary choice of
exchange. On the global class $z\otimes1$, replacing this
crossing by the ordinary swap while keeping the product fixed
gives the nonzero residual $2\delta_{12}$ of
Example~\ref{gh:crossing-example}. The coefficient of the
diagonal therefore measures a change in the actual operation.
In particular, Theorem~\ref{sh:hopf} does not make the full Hall
carrier an ordinary-symmetry bialgebra.
\end{example}

\subsection{Comparison on the cyclic companion family}

Trivialize the common $E$-torsor and work on the support atlas.
Here $c_i:\mathcal M_i\to\operatorname{Sym}^i(E)$ denotes
the unquotiented support map. In the coordinate disk $D$, use
$Y=\operatorname{Sym}^3(D)$, $X=D^3$, $m:X\to Y$, and
$j:U\hookrightarrow X$ from Section~\ref{sec:companion-family}.
For the Springer comparison, label the ordered endpoints
subobject first. This reverses the quotient-first labels in
\eqref{cf:endpoint-lines} without changing the finite map $m$.
The cyclic module
\[
 \mathcal O_Y[\zeta]/
       (\zeta^3-e_1\zeta^2+e_2\zeta-e_3)
\]
defines a section $f:Y\to\mathcal M_3$ of $c_3$.
Proposition~\ref{cf:families} identifies the actual full-flag
pullback with $m$ and the strict full-splitting pullback with
$mj$. Their coefficient objects and restriction map are
\[
 \mathcal F=Rm_*\mathbb Q_X,
 \qquad \mathcal G=R(mj)_*\mathbb Q_U,
 \qquad r:\mathcal F\longrightarrow\mathcal G.
\]
Write $u_{\mathcal F}:\mathbb Q_Y\to\mathcal F$ for the unit
and $t_{\mathcal F}:\mathcal F\to\mathbb Q_Y$ for the proper
trace. Its weights are $1,2,6$ on the three multiplicity types,
as computed in Proposition~\ref{cf:proper-coefficient}, and
$t_{\mathcal F}u_{\mathcal F}=6\operatorname{id}$.

\begin{theorem}[The support-family Hall comparison]
\label{gh:support-family}
There are canonical maps
\[
 e_{\mathcal F}:\mathcal V_3|_Y\longrightarrow\mathcal F,
 \qquad
 e_{\mathcal G}:\mathcal V_3|_Y\longrightarrow\mathcal G,
 \qquad
 \varepsilon:\mathcal B_3|_Y\longrightarrow\mathbb Q_Y
\]
such that
\begin{equation}\label{sh:family-identities}
 \begin{gathered}
 e_{\mathcal G}=r e_{\mathcal F},\qquad
 e_{\mathcal F}j_3=u_{\mathcal F}\varepsilon,\\
 \varepsilon P_3=t_{\mathcal F}e_{\mathcal F}.
 \end{gathered}
\end{equation}
The first evaluation intertwines the supported Springer action
with the permutation action on the proper finite flag family.
The equations commute with both three-step parenthesizations,
including the nonsplit cyclic fibre and its endpoint lines.
\end{theorem}
\begin{proof}
The first map is counit evaluation followed by proper base change:
\[
 \mathcal V_3|_Y
 \xrightarrow{\ \gamma_3\ }
 (Rc_{3,*}\mathrm{Spr}_3)|_Y
 \longrightarrow f^*\mathrm{Spr}_3
 \xrightarrow{\ \sim\ }\mathcal F.
\]
The middle arrow applies $f^*$ to
$c_3^*Rc_{3,*}\mathrm{Spr}_3\to\mathrm{Spr}_3$, using
$c_3f=\operatorname{id}_Y$. This construction is natural for
the Springer action. On $Y$, its target action is the deck
action extended to the proper finite map $m$. Indeed, the
pullback action has that restriction on distinct supports, and
$\mathcal F[3]$ is their finite-map intermediate extension;
uniqueness identifies the two actions.

For the full direct-sum map
$s:(\mathcal M_1)^3\to\mathcal M_3$, define the second
evaluation by
\[
 \mathcal V_3|_Y=R(c_3s)_*\mathbb Q|_Y
 \longrightarrow f^*Rs_*\mathbb Q
 \longrightarrow\mathcal G.
\]
The first arrow is counit evaluation for $Rs_*\mathbb Q$;
the second is the ordinary base-change transformation to the
strict splitting pullback. No inverse to that transformation
is assumed. On the distinct-support open $V\subset Y$, an
ordered decomposition and its full flag determine each other,
so $e_{\mathcal G}$ and $r e_{\mathcal F}$ agree there.
But $\mathcal G=R(j_V)_*\mathcal L$ for the local system
$\mathcal L=(m|_U)_*\mathbb Q_U$ on $V$. Adjunction gives
\[
 \operatorname{Map}(\mathcal V_3|_Y,\mathcal G)
 \simeq
 \operatorname{Map}(j_V^*\mathcal V_3,\mathcal L).
\]
It follows that their open identification extends to the
derived equality $e_{\mathcal G}=r e_{\mathcal F}$.

Define $\varepsilon$ by counit evaluation of
$Rc_{3,*}\mathbb Q$ along $f$.
Unit--counit naturality and \eqref{sh:gamma} give
$e_{\mathcal F}j_3=u_{\mathcal F}\varepsilon$.
The equivariance just proved and the full symmetrizer give
\[
 \begin{aligned}
 u_{\mathcal F}\varepsilon P_3
  &=e_{\mathcal F}j_3P_3
    =\sum_{w\in S_3}w\,e_{\mathcal F}\\
  &=u_{\mathcal F}t_{\mathcal F}e_{\mathcal F}.
 \end{aligned}
\]
Here $u_{\mathcal F}t_{\mathcal F}=\sum_w w$ is the finite
proper trace identity, with the actual ramification weights.
The left inverse $t_{\mathcal F}/6$ of $u_{\mathcal F}$ permits
cancellation and proves the last equation of
\eqref{sh:family-identities}. Thus the trace comparison does
not assume that arbitrary pullback preserves a Verdier trace.

The two three-step decompositions are refinements of these full
direct-sum and full-flag diagrams. Their supported comparisons
are those of Proposition~\ref{sh:block-action} and
Theorem~\ref{sh:operation-comparison}. Proper base change,
the displayed counits, and the open-restriction adjunction
commute with these refinements. Their quotient lines are the
actual endpoint lines of \eqref{cf:endpoint-lines}, with the
specified reversal when written subobject first.
Consequently the comparisons retain both parenthesizations on
the whole cyclic family.
\end{proof}

On the separate labelled base $T$, ordinary base change of $\mathcal G$
is still the map $v\mapsto(\sigma^*v)_\sigma$ into the six
strict-splitting copies. The proper pullback retains its
collision-incidence kernel. The maps and cone ranks of
Section~\ref{sec:companion-family} are therefore unchanged.
Theorem~\ref{gh:support-family} supplies arrows between
$\mathcal V_3|_Y$, $\mathcal F$, and $\mathcal G$; it does not
identify these objects. In particular, it supplies no trace
$\mathcal G\to\mathbb Q_Y$ extending the generic proper trace,
which is excluded by Theorem~\ref{cf:no-trace}.

The endpoint twists also give a direct obstruction to a
multiplicative cyclic evaluation. To state it globally, let
$\mathcal D_n$ be the universal degree-$n$ divisor on the
universal elliptic torsor over $Y_n$. A descended invertible
module $N$ on $\mathcal D_n$ defines a cyclic section
$s_n^N:Y_n\to\overline{\mathcal M}_n$. Its adjunction unit
gives the evaluation
\[
 \varepsilon_n^N:\mathcal B_n
 \longrightarrow R\overline c_{n,*}Rs_{n,*}^N\mathbb Q
 \simeq\mathbb Q_{Y_n}.
\]
It sends $1$ to $1$. On the coordinate family at charge three
and for the trivial module, this is the map $\varepsilon$
in \eqref{sh:family-identities}.
Write $\mu^{\mathrm{quot}}=\mu\tau$ for the ordinary endpoint
reversal described after Theorem~\ref{gh:global-hopf}.

\begin{proposition}\label{sh:cyclic-defect}
Let $z$ be the charge-one stabilizer Chern class and let
$\mathfrak D_2\subset Y_2$ be the reduced discriminant divisor.
For the quotient-first Hall product,
\begin{equation}\label{sh:binary-cyclic-defect}
 \varepsilon_2^N\mu^{\mathrm{quot}}_{1,1}
             (1\boxtimes z-z\boxtimes1)
       =-[\mathfrak D_2]\ne0.
\end{equation}
No family of these constant Betti evaluations is multiplicative
for a bilinear multiplication on the evaluated recipients.
Nonzero scalar rescalings of the evaluations do not remove the
obstruction.
\end{proposition}
\begin{proof}
Use the quotient-first endpoint order of
\eqref{cf:endpoint-lines}. In charge two the lines are
$e_1^*N$ and $e_2^*N\otimes\mathcal O(-\Delta_{12})$ on
$\widetilde Y_2$. This follows from the same ideal sequence
for a monic factorization used in
Proposition~\ref{cf:families}; it retains the nonreduced flag
fibre. With $\ell_i=c_1(e_i^*N)$, the proper cyclic trace of
the displayed insertion is
\[
 (q_2)_*(\ell_2-\ell_1-[\Delta_{12}])
          =-(q_2)_*[\Delta_{12}].
\]
The two $\ell_i$ have equal pushforward, since the transposition
exchanges them. The map $q_2$ takes its ramification divisor
with degree one to the reduced discriminant, so the last
expression is $-[\mathfrak D_2]$.
This is the proper trace of the actual cyclic flag pullback.
Equivalently, the cyclic section is a smooth atlas of the cyclic
locus: its trivializations form a torsor under the unit group
of the rank-two finite support algebra, a smooth open subset of
that algebra. Smooth base change therefore gives the same trace
comparison.

The complete pair map
$\iota:E\to Y_2$, $\delta\mapsto[0]+[\delta]$, satisfies
$\iota^*\mathfrak D_2=2[0]$. Locally the polynomial is
$x(x-t)$ and its discriminant is $t^2$. Hence the pullback of
\eqref{sh:binary-cyclic-defect} has integral $-2$, proving
nonvanishing. But $Y_1=[E/E]=\mathrm{pt}$, so every such
evaluation sends $z$ to zero. A bilinear multiplication after
evaluation would therefore kill both terms of the insertion,
contradicting \eqref{sh:binary-cyclic-defect}. Rescaling by
nonzero constants preserves this contradiction.
\end{proof}

The subobject-first convention reverses the sign of the
antisymmetric insertion; its nonvanishing is the same.

The construction proves the supported crossing and Hopf comparison
for the normalized rational spherical Betti coefficients.
Interacting critical coefficients require their own orientation,
coefficient, and gluing comparisons. Their possible negative
cohomological degrees also lie outside the nonnegative
fibre-quotient argument. An independently specified
field or Ran-space realization, compact primitive recognition
beyond the ordinary even quotient, and physical state
reconstruction also require their actual comparison maps.
The compact realization questions in
Section~\ref{sec:compact-realization} retain those hypotheses.
